% ============================================================================
% Fichier  : partie2/P02_C07_methodologies_standards.tex
% Titre    : Méthodologies et Standards Internationaux
% Auteur   : MINKA MI NGUIDJOÏ Thierry Emmanuel
% Version  : 1.0 -- Juin 2026
% Statut   : RÉDIGÉ
% Sources  : 7_methodologies_investigation.tex, 5_cadre_normatif_global.tex,
%            M4_Cadre_Normatif_International_PNC.docx, processus_investigatif_universel.pdf
% ============================================================================

\chapter{Méthodologies et Standards Internationaux}
\label{chap:meth-std}

\epigraph{%
    « Digital forensics is not just about recovering data, it's about
    understanding the context in which the data existed. »%
}{--- Joshua I. James}

\niveauChapitre{Fondamental}{%
    Lecture du Chapitre~\ref{chap:PIU} (PIU) recommandée~: ce chapitre
    positionne les standards internationaux par rapport au PIU et démontre
    leur cohérence.%
}

\begin{boxobjectifs}
\begin{itemize}
    \item Maîtriser les cinq standards forensiques mondiaux de référence~:
          DFRWS, ISO/IEC~27037/41/42/43, NIST~SP~800-86, RFC~3227, ACPO.
    \item Comprendre la logique des standards~: pourquoi ils existent,
          ce qu'ils imposent, et ce qu'ils laissent libre.
    \item Mettre en correspondance chaque standard avec les phases du PIU
          pour montrer leur cohérence structurelle.
    \item Appliquer le framework d'analyse \termecle{DICES} pour comparer
          les méthodologies mondiales.
    \item Identifier les équivalences entre standards internationaux et
          droit camerounais (\loicam{2010/012}, \loicam{2024/017}).
    \item Choisir le standard adapté à une investigation selon son contexte.
\end{itemize}
\end{boxobjectifs}

\bigskip

% ============================================================================
\section*{Standards et PIU~: une relation de complémentarité}
% ============================================================================

Le Processus Investigatif Universel (PIU) présenté au Chapitre~\ref{chap:PIU}
est un cadre général adaptatif. Les standards internationaux qui font l'objet
de ce chapitre sont ses fondements normatifs~: ils précisent \emph{comment}
chaque phase doit être conduite pour être reconnue par les juridictions et
institutions du monde entier. La relation n'est pas de concurrence mais de
complémentarité~: le PIU structure, les standards légitiment.

\separateur

% ============================================================================
\section{Architecture des Standards Forensiques~: vue d'ensemble}
\label{sec:std-architecture}
% ============================================================================

\subsection{Pourquoi des standards~?}

Sans standards communs, deux investigateurs de deux pays différents peuvent
analyser la même preuve avec des méthodes incompatibles et produire des
conclusions qui se contredisent sans qu'aucune expertise externe ne puisse
trancher. Les standards résolvent trois problèmes fondamentaux~:
\begin{enumerate}
    \item \textbf{Interopérabilité}~: les preuves collectées selon ISO~27037
          au Cameroun sont recevables dans un tribunal belge ou canadien~;
    \item \textbf{Opposabilité}~: un expert qui documente sa conformité à
          NIST~SP~800-86 peut répondre à n'importe quelle attaque méthodologique~;
    \item \textbf{Formation}~: les standards constituent une base d'apprentissage
          commune pour les investigateurs du monde entier.
\end{enumerate}

\subsection{Cartographie des standards}

\begin{tableaucours}{p{3.0cm}p{2.5cm}p{2.5cm}p{4.0cm}}{%
    \tableauHeader{Standard} &
    \tableauHeader{Organisme} &
    \tableauHeader{Domaine} &
    \tableauHeader{Contribution principale}
}
DFRWS 2001 & DFRWS & Processus global &
    Premier modèle formel~: 6~phases \\
\norme{ISO/IEC 27037} & ISO/IEC & Collecte & Rôles DEFR/DES,
    identification et acquisition \\
\norme{ISO/IEC 27041} & ISO/IEC & Méthodes &
    Validation des méthodes d'investigation \\
\norme{ISO/IEC 27042} & ISO/IEC & Analyse &
    Analyse et interprétation des preuves \\
\norme{ISO/IEC 27043} & ISO/IEC & Processus complet &
    Principes et processus intégrés \\
\norme{NIST SP 800-86} & NIST (USA) & IR + Forensique &
    Intégration réponse aux incidents \\
\norme{RFC 3227} & IETF & Volatilité &
    Ordre de collecte, bonnes pratiques \\
ACPO Guide & ACPO (UK) & Déontologie &
    4~principes, responsabilité légale \\
SANS FOR508 & SANS Institute & IR avancée &
    Forensique entreprise, threat hunting \\
\end{tableaucours}
\vspace{-0.8em}
\begin{center}{\small\itshape Tableau~7.1 --- Cartographie des standards forensiques}
\end{center}

\subsection{Correspondance standards / phases PIU}

\begin{tableaucours}{p{4.2cm}p{2.0cm}p{2.0cm}p{2.0cm}p{2.0cm}}{%
    \tableauHeader{Standard} &
    \tableauHeader{Ph.1\\Init.} &
    \tableauHeader{Ph.2\\Acquis.} &
    \tableauHeader{Ph.3\\Analyt.} &
    \tableauHeader{Ph.4--5\\Form.}
}
\norme{ISO/IEC 27037} & $\bullet$ & $\bullet\bullet\bullet$ & & \\
\norme{ISO/IEC 27041} & & $\bullet$ & $\bullet\bullet$ & \\
\norme{ISO/IEC 27042} & & & $\bullet\bullet\bullet$ & \\
\norme{ISO/IEC 27043} & $\bullet$ & $\bullet$ & $\bullet$ & $\bullet$ \\
\norme{NIST SP 800-86} & $\bullet$ & $\bullet\bullet$ & $\bullet\bullet$ & $\bullet$ \\
\norme{RFC 3227} & & $\bullet\bullet\bullet$ & & \\
ACPO Guide & $\bullet$ & $\bullet\bullet$ & $\bullet$ & $\bullet$ \\
SANS FOR508 & $\bullet\bullet$ & $\bullet$ & $\bullet\bullet$ & $\bullet$ \\
\end{tableaucours}
\vspace{-0.8em}
\begin{center}{\small\itshape Tableau~7.2 --- Correspondance standards / phases PIU~:
$\bullet$ pertinent~; $\bullet\bullet$ central~; $\bullet\bullet\bullet$ définitoire}
\end{center}

% ============================================================================
\section{DFRWS~: la Fondation Conceptuelle (2001)}
\label{sec:dfrws}
% ============================================================================

Le \termecle{Digital Forensics Research Workshop}\index{DFRWS} publie en 2001
le premier cadre formel de forensique numérique. Son importance est
fondatrice~: il établit un vocabulaire commun et une séquence logique que
tous les standards suivants reprendront et affineront.

\subsection{Les six phases DFRWS}

\begin{figure}[H]
\centering
\begin{tikzpicture}[
    box/.style={draw=CouleurPrimaire, fill=CouleurDef,
        rounded corners=3pt, minimum width=2.5cm, minimum height=0.9cm,
        font=\small\bfseries\color{CouleurPrimaire}, align=center},
    fl/.style={-{Stealth[length=4pt]}, thick, color=CouleurSecondaire}
]
\node[box] (id)   at (0,0)    {Identifi-\\cation};
\node[box] (pr)   at (2.7,0)  {Préser-\\vation};
\node[box] (co)   at (5.4,0)  {Collec-\\tion};
\node[box] (ex)   at (8.1,0)  {Examina-\\tion};
\node[box] (an)   at (5.4,-1.8) {Analyse};
\node[box] (pe)   at (2.7,-1.8) {Présen-\\tation};
\draw[fl] (id) -- (pr);
\draw[fl] (pr) -- (co);
\draw[fl] (co) -- (ex);
\draw[fl] (ex) -- (an);
\draw[fl] (an) -- (pe);
% Correspondance PIU
\node[font=\tiny\itshape, color=CouleurSecondaire] at (0,-0.9) {PIU~E1--E3};
\node[font=\tiny\itshape, color=CouleurSecondaire] at (2.7,-0.9) {PIU~E4--E5};
\node[font=\tiny\itshape, color=CouleurSecondaire] at (5.4,-0.9) {PIU~E4};
\node[font=\tiny\itshape, color=CouleurSecondaire] at (8.1,-0.9) {PIU~E7--E9};
\node[font=\tiny\itshape, color=CouleurSecondaire] at (5.4,-2.7) {PIU~E9--E10};
\node[font=\tiny\itshape, color=CouleurSecondaire] at (2.7,-2.7) {PIU~E12--E13};
\end{tikzpicture}
\caption{Les 6~phases DFRWS et leurs correspondances avec le PIU}
\label{fig:dfrws}
\end{figure}

\begin{boxdef}
\textbf{Identification}~: reconnaître l'existence d'un incident numérique
et en définir le périmètre. \textbf{Préservation}~: protéger les preuves
de toute altération --- phase la plus critique, car une erreur ici est
irréparable. \textbf{Collection}~: acquérir les preuves selon des méthodes
documentées. \textbf{Examination}~: extraire les informations pertinentes.
\textbf{Analyse}~: corréler et reconstruire les événements.
\textbf{Presentation}~: communiquer les conclusions de façon défendable.
\end{boxdef}

\subsection{Limite historique du DFRWS et apport du PIU}

Le modèle DFRWS est \textbf{linéaire}. En pratique, la Phase~3
(\textit{Collection}) et la Phase~4 (\textit{Examination}) s'entrelacent~:
chaque artefact découvert lors de l'examen peut déclencher une nouvelle
collecte. Le PIU formalise cette réalité avec la boucle itérative
E4--E6~$\leftrightarrow$~E7--E10. Le DFRWS reste la référence conceptuelle~;
le PIU est l'implémentation opérationnelle.

% ============================================================================
\section{La Suite ISO/IEC 2703x~: le Standard International de Référence}
\label{sec:iso-2703x}
% ============================================================================

\subsection{Architecture de la famille}

La famille ISO/IEC~2703x constitue l'architecture normative la plus complète
et la plus internationalement reconnue pour la forensique numérique. Ses
quatre normes couvrent l'intégralité du cycle de vie des preuves numériques.

% ---- ISO 27037 ----
\subsection{ISO/IEC~27037:2012 --- Identification, Collecte et Préservation}

\norme{ISO/IEC~27037:2012} --- \textit{Guidelines for identification,
collection, acquisition and preservation of digital evidence} ---
est \textbf{le standard le plus important pour le travail de terrain}.
Son apport central est la définition de deux rôles distincts et de quatre
principes fondamentaux.

\subsubsection*{Les deux rôles}

\begin{itemize}
    \item \textbf{DEFR} (\textit{Digital Evidence First Responder})~:
          le premier intervenant sur la scène. Son rôle est la préservation
          initiale. Son erreur est irréparable. C'est l'OPJ qui arrive le
          premier dans l'affaire MBONGO.
    \item \textbf{DES} (\textit{Digital Evidence Specialist})~:
          l'analyste en laboratoire. Il reçoit les preuves préservées par
          le DEFR et conduit l'examination et l'analysis.
\end{itemize}

\subsubsection*{Les quatre principes ISO~27037}

\begin{enumerate}
    \item \textbf{Pertinence}~: la collecte doit être ciblée et justifiée.
          Ne collecter que ce qui est pertinent pour l'investigation.
    \item \textbf{Fiabilité}~: les méthodes utilisées doivent être
          reproductibles et vérifiables.
    \item \textbf{Suffisance}~: la collecte doit être exhaustive dans
          le périmètre défini. Pas de lacunes qui permettraient de contester
          la représentativité de la base probatoire.
    \item \textbf{Documentation}~: traçabilité complète de chaque action.
\end{enumerate}

\subsubsection*{Procédure de saisie selon ISO~27037}

\begin{lstlisting}[language=bash_forensique,
    caption={Checklist de saisie ISO~27037 --- séquence minimale}]
# Étape 1 — Identification préliminaire
# Type de dispositif, état (allumé/éteint), connexions actives
# -> Appliquer l'ordre de volatilité RFC 3227 (cf. section suivante)

# Étape 2 — Documentation photographique
# Photos du contexte avant toute manipulation
# Horodatage visible (montre + journal de bord)

# Étape 3 — Isolation du dispositif
# Mobile : mode avion ou sac de Faraday
# PC allumé : capturer RAM avant isolation réseau

# Étape 4 — Acquisition (write-blocker obligatoire)
# FTK Imager : File -> Create Disk Image -> Format E01
# Vérifier que le write-blocker est actif avant montage

# Étape 5 — Vérification d'intégrité (hash SHA-256 minimum)
sha256sum /dev/sdb > hash_source.txt  # Linux
certutil -hashfile image.E01 SHA256   # Windows
# Conserver les deux hash dans le formulaire custody

# Étape 6 — Scellement et transport
# Sachet inviolable, étiquette, formulaire custody signé
\end{lstlisting}

\begin{boxjuris}
\textbf{Compatibilité ISO/IEC~27037 et droit camerounais~:}

Les articles~52--59 de la \loicam{2010/012} imposent~: présence des parties,
copie réalisée en leur présence, documentation. Ces exigences sont
\textbf{substantiellement alignées} sur ISO/IEC~27037. Un investigateur
respectant le standard international respecte de facto le droit camerounais.
Inversement, une investigation conforme aux articles~52--59 satisfait
aux critères DEFR d'ISO/IEC~27037.
\end{boxjuris}

% ---- ISO 27041 ----
\subsection{ISO/IEC~27041:2015 --- Validation des Méthodes}

\norme{ISO/IEC~27041:2015} --- \textit{Guidance on assuring suitability and
adequacy of incident investigative method} --- répond à la question~:
\emph{comment s'assurer qu'une méthode d'investigation est scientifiquement
valide~?} Ce standard est particulièrement important pour l'opposabilité~:
il fournit le cadre pour justifier le choix d'un outil ou d'une technique
devant un tribunal.

Quatre domaines de méthodes validées~:
\begin{itemize}
    \item \textbf{Live Forensics}~: investigation sur système en cours
          d'exécution~; méthode conforme à RFC~3227.
    \item \textbf{Dead Forensics}~: investigation sur système éteint~;
          méthode conforme à NIST~SP~800-86.
    \item \textbf{Network Forensics}~: capture et analyse de trafic~;
          standards IETF (RFC~3227).
    \item \textbf{Mobile Forensics}~: quatre niveaux d'acquisition~;
          NIST~SP~800-101r1.
\end{itemize}

% ---- ISO 27042 ----
\subsection{ISO/IEC~27042:2015 --- Analyse et Interprétation}

\norme{ISO/IEC~27042:2015} --- \textit{Guidelines for the analysis and
interpretation of digital evidence} --- est le standard de la Phase~3 du PIU.
Son apport clé est la formalisation de la séquence d'analyse~:

\begin{enumerate}
    \item \textbf{Préparation}~: environnement d'analyse isolé, outils
          validés et versionnés.
    \item \textbf{Extraction}~: récupération des données depuis les
          images forensiques.
    \item \textbf{Analyse}~: application des techniques appropriées.
    \item \textbf{Interprétation}~: contextualisation~; c'est ici que
          la distinction constatation/interprétation s'applique.
    \item \textbf{Rapport}~: documentation rigoureuse des découvertes,
          des méthodes utilisées et des limitations.
\end{enumerate}

ISO~27042 renforce ce que le PIU appelle E9bis (falsification active)~:
le standard impose de \textbf{tester les conclusions alternatives avant
de les rejeter}, pas simplement de choisir l'explication la plus probable.

% ---- ISO 27043 ----
\subsection{ISO/IEC~27043:2015 --- Principes et Processus Intégrés}

\norme{ISO/IEC~27043:2015} est le standard de niveau le plus élevé~:
il intègre les trois précédents dans un processus complet incluant la
\termecle{readiness} (préparation pré-incident) et la gestion
post-investigation. Sa représentation canonique~:

\begin{figure}[H]
\centering
\begin{tikzpicture}[
    ph/.style={draw=CouleurSecondaire, fill=CouleurDef, rounded corners=3pt,
        font=\footnotesize\bfseries\color{CouleurPrimaire}, align=center,
        minimum height=0.8cm, inner xsep=4pt},
    fl/.style={-{Stealth[length=4pt]}, thick, color=CouleurSecondaire}
]
\node[ph, minimum width=2.0cm] (r) at (0,0) {Readiness};
\node[ph, minimum width=1.8cm] (d) at (2.2,0) {Detection};
\node[ph, minimum width=2.0cm] (ir) at (4.3,0) {Initial\\Response};
\node[ph, minimum width=1.8cm] (st) at (6.3,0) {Strategy};
\node[ph, minimum width=1.8cm] (c) at (0,-1.5) {Collection};
\node[ph, minimum width=1.8cm] (an) at (2.2,-1.5) {Analysis};
\node[ph, minimum width=2.0cm] (pr) at (4.3,-1.5) {Presentation};
\node[ph, minimum width=2.4cm] (pi) at (6.3,-1.5) {Post-Investigation};
\draw[fl] (r) -- (d);
\draw[fl] (d) -- (ir);
\draw[fl] (ir) -- (st);
\draw[fl] (st) -- (c);
\draw[fl] (c) -- (an);
\draw[fl] (an) -- (pr);
\draw[fl] (pr) -- (pi);
% Retours possibles
\draw[fl, dashed, color=CouleurAccent]
    (an.south) to[bend right=20]
    node[below, font=\tiny\itshape, color=CouleurAccent] {besoin supplémentaire}
    (c.south);
\end{tikzpicture}
\caption{Processus ISO/IEC~27043~: 8~phases avec boucle retour}
\label{fig:iso27043}
\end{figure}

% ============================================================================
\section{NIST SP~800-86~: Forensique et Réponse aux Incidents}
\label{sec:nist-800-86}
% ============================================================================

\norme{NIST SP~800-86} --- \textit{Guide to Integrating Forensic Techniques
into Incident Response} --- est la référence américaine. Sa contribution
distinctive par rapport aux normes ISO est l'\textbf{intégration de la
forensique dans la réponse aux incidents}~: le NIST ne traite pas la
forensique comme une activité autonome mais comme une composante de la
gestion des incidents de sécurité.

\subsection{Les quatre phases NIST}

\begin{tableaucours}{p{2.8cm}p{4.5cm}p{4.0cm}}{%
    \tableauHeader{Phase NIST} &
    \tableauHeader{Activités} &
    \tableauHeader{Correspondance PIU}
}
Collection & Priorisation des données~; préservation~; chaîne de custody &
    E4--E5 (Acquisition) \\
Examination & Extraction des données~; revue manuelle~; analyse automatisée &
    E7--E8 (Investigation) \\
Analysis & Reconstruction de la chronologie~; corrélation~; attribution &
    E9--E10 (Investigation analytique) \\
Reporting & Synthèse exécutive~; détails techniques~; recommandations &
    E11--E12 (Formalisation) \\
\end{tableaucours}
\vspace{-0.8em}
\begin{center}{\small\itshape Tableau~7.3 --- Phases NIST SP~800-86
et correspondances PIU}\end{center}

\subsection{Distinction Examination / Analysis~: règle d'or du rapport}

La contribution épistémologique la plus importante du NIST~SP~800-86 est la
séparation stricte entre~:
\begin{itemize}
    \item \textbf{Examination}~: documenter ce qu'on \emph{trouve}.
          Résultats factuels, vérifiables, non ambigus. \emph{Exemple~:}
          « Le fichier \texttt{enc\_all.ps1} est présent dans
          \texttt{\%TEMP\%} avec un horodatage de création le
          12/03/2026~à~02h17. »
    \item \textbf{Analysis}~: interpréter ce que cela \emph{signifie}.
          Conclusions raisonnées, avec niveau de certitude explicite.
          \emph{Exemple~:} « La présence de ce script à cet horodatage,
          corrélée avec la connexion RDP de 02h14, indique \emph{avec
          une haute confiance} que le script a été déposé par l'attaquant
          lors de sa connexion. »
\end{itemize}

\begin{boxalert}
\textbf{Un avocat peut attaquer une interprétation~; il ne peut pas attaquer
un hash qui correspond.} Séparer strictement les deux types d'énoncés dans
le rapport est la protection la plus efficace contre les attaques méthodologiques
en contre-interrogatoire.
\end{boxalert}

% ============================================================================
\section{RFC~3227~: l'Ordre de Volatilité}
\label{sec:rfc3227}
% ============================================================================

La RFC~3227 (\textit{Guidelines for Evidence Collection and Archiving},
Brezinski \& Killalea, 2002) est la référence technique pour l'ordre de
collecte sur les systèmes en cours d'exécution. Elle complète ISO~27037
en précisant la \textbf{priorité de collecte selon la volatilité des données}.

\subsection{L'ordre de volatilité canonique}

\begin{tableaucours}{p{0.8cm}p{4.5cm}p{6.0cm}}{%
    \tableauHeader{Rang} &
    \tableauHeader{Type de données} &
    \tableauHeader{Durée de vie et outil}
}
1 & Registres CPU, cache &
    Nanosecondes~; non capturable en pratique \\
2 & Mémoire système (RAM) &
    Jusqu'à extinction~; FTK~Imager Portable (\texttt{Capture Memory}) \\
3 & État réseau (tables de routage, ARP, connexions) &
    Secondes à minutes~; \texttt{netstat -ano}, \texttt{arp -a} \\
4 & Processus en cours d'exécution &
    Minutes~; \texttt{tasklist} (Windows) ou \texttt{ps aux} (Linux) \\
5 & Fichiers temporaires, logs actifs &
    Heures~; \texttt{dir /T} ou \texttt{ls -lt /tmp} \\
6 & Disque dur (système de fichiers) &
    Mois à années~; FTK~Imager (image E01 avec hash) \\
7 & Logs système distants &
    Jours à semaines selon rotation~; réquisition opérateur \\
8 & Configuration physique, topologie réseau &
    Stables~; documentés en dernier \\
\end{tableaucours}
\vspace{-0.8em}
\begin{center}{\small\itshape Tableau~7.4 --- Ordre de volatilité RFC~3227}
avec outils de collecte}\end{center}

\begin{boxPNC}
\textbf{Application OPJ~: règle mnémotechnique RFC~3227}

\textbf{Sur un système allumé~: RAM avant disque, réseau avant fichiers.}

Arrivée sur la scène~: (1)~capturer la mémoire RAM en premier
(\volatilite{2}{RAM})~; (2)~documenter les connexions réseau
(\volatilite{3}{Réseau})~; (3)~lister les processus actifs
(\volatilite{4}{Processus})~; (4)~éteindre proprement~;
(5)~acquérir le disque avec write blocker
(\volatilite{6}{Disque}).

Sur un système éteint~: passer directement à l'étape~(5).
Ne jamais allumer un système éteint sans write blocker actif.
\end{boxPNC}

% ============================================================================
\section{Les Quatre Principes ACPO~: Responsabilité et Déontologie}
\label{sec:acpo}
% ============================================================================

Le guide ACPO (\textit{Association of Chief Police Officers}, Royaume-Uni,
2012) est la référence déontologique mondiale pour les forces de l'ordre.
Son importance tient à sa simplicité~: quatre principes mémorables qui
encadrent toute l'activité forensique policière.

\begin{tableaucours}{p{2.2cm}p{9.5cm}}{%
    \tableauHeader{Principe} &
    \tableauHeader{Formulation et implications}
}
\textbf{Principe~1} & \textbf{Aucune action ne doit modifier les données.}
    Toute opération de collecte doit laisser les données originales intactes.
    Fondement du write blocker, du calcul de hash et du travail sur copies. \\
\textbf{Principe~2} & \textbf{Si une modification est nécessaire,
    elle requiert la compétence technique et doit être documentée.}
    Un OPJ non formé ne doit pas procéder à l'acquisition~: il appellera
    un spécialiste. Toute modification est explicitement justifiée. \\
\textbf{Principe~3} & \textbf{Un audit trail complet doit être maintenu.}
    Chaque action, chaque accès, chaque transfert est documenté. C'est la
    traduction procédurale de la chaîne de custody. \\
\textbf{Principe~4} & \textbf{Le responsable de l'investigation est
    responsable de la conformité.} L'investigateur en chef endosse personnellement
    la responsabilité du respect de ces principes pour l'ensemble de l'équipe. \\
\end{tableaucours}
\vspace{-0.8em}
\begin{center}{\small\itshape Tableau~7.5 --- Les quatre principes ACPO}
\end{center}

\begin{boxCRO}
\textbf{Analyse CRO --- Principes ACPO}

Les principes ACPO maximisent $\mathcal{R}$ et $\mathcal{O}$ au détriment
partiel de la flexibilité opérationnelle~: ils imposent une rigueur absolue
même dans les situations d'urgence. Cette tension est délibérée~: dans
le contexte policier, la $\mathcal{O}$ doit primer pour que les preuves
collectées mènent à des condamnations incontestables.

Le Principe~4 (responsabilité personnelle) est la réponse institutionnelle
au problème d'attribution~: quand la chaîne est violée, il doit être possible
d'identifier précisément qui est responsable.

\cro{$\sim$}{$\uparrow\uparrow$}{$\uparrow\uparrow$}
\end{boxCRO}

% ============================================================================
\section{SANS FOR508~: Forensique d'Entreprise et Threat Hunting}
\label{sec:sans-for508}
% ============================================================================

Le cours \norme{SANS FOR508} (\textit{Advanced Incident Response, Threat Hunting,
and Digital Forensics}) représente l'état de l'art de la forensique
d'entreprise. Contrairement aux normes ISO qui s'appliquent à tous les
contextes, FOR508 cible spécifiquement les investigations à grande échelle
sur des réseaux d'entreprise compromis.

\subsection{La méthodologie en six phases}

\begin{enumerate}
    \item \textbf{Preparation}~: plan de réponse aux incidents, validation des
          outils, formation de l'équipe~;
    \item \textbf{Identification}~: développement des IoC
          (\textit{Indicators of Compromise}), intégration de la Threat
          Intelligence, détection d'anomalies~;
    \item \textbf{Containment}~: isolation à court terme, planification
          de l'éradication, préservation des preuves~;
    \item \textbf{Eradication}~: suppression du malware, correction des
          vulnérabilités, durcissement des systèmes~;
    \item \textbf{Recovery}~: restauration des systèmes, surveillance
          renforcée, tests de validation~;
    \item \textbf{Lessons Learned}~: revue post-incident, amélioration des
          processus, mise à jour de la documentation.
\end{enumerate}

\subsection{Apports distinctifs de FOR508 en 2025}

\begin{itemize}
    \item \textbf{Forensique mémoire avec Volatility~3}~: syntaxe mise à
          jour (\texttt{python3 vol.py -f mem.dmp windows.pslist}),
          nouveaux plugins pour artefacts modernes.
    \item \textbf{EDR Forensics}~: exploitation des journaux produits par les
          outils de détection/réponse sur les endpoints (CrowdStrike, SentinelOne)~;
          souvent plus complets que les logs Windows natifs.
    \item \textbf{Enterprise-scale IR}~: investigation simultanée sur
          des centaines de machines avec des outils comme Velociraptor ou KAPE.
    \item \textbf{Threat Hunting}~: recherche proactive de compromissions
          latentes \emph{avant} qu'une alerte soit déclenchée.
\end{itemize}

\begin{boiteRemarque}{SANS FOR508 vs ISO~27037~: quelle différence~?}
ISO~27037 est un standard \emph{normatif}~: il définit ce que vous devez faire.
SANS FOR508 est un cours \emph{opérationnel}~: il vous apprend à le faire.
FOR508 va beaucoup plus loin dans les techniques spécifiques~; ISO~27037 a
une portée légale internationale. En pratique, un bon investigateur maîtrise
les deux~: la légitimité normative (ISO) et les techniques avancées (SANS).
\end{boiteRemarque}

% ============================================================================
\section{Analyse Comparative~: le Framework DICES Appliqué}
\label{sec:comparaison-dices}
% ============================================================================

\subsection{Grille comparative DICES des standards}

Chaque standard peut être analysé selon les cinq axes DICES
(\textbf{D}octrine, \textbf{I}nfrastructure, \textbf{C}apacités,
\textbf{E}cosystème, \textbf{S}tratégie)~:

\begin{tableaucours}{p{2.8cm}p{2.2cm}p{2.0cm}p{2.2cm}p{2.5cm}}{%
    \tableauHeader{Standard} &
    \tableauHeader{Doctrine} &
    \tableauHeader{Public cible} &
    \tableauHeader{Portée légale} &
    \tableauHeader{Spécialité forensique}
}
\norme{ISO/IEC 27037} & Collecte terrain & DEFR / DES &
    Internationale & Acquisition et préservation \\
\norme{ISO/IEC 27042} & Analyse & DES & Internationale & Analyse et interprétation \\
\norme{ISO/IEC 27043} & Processus global & Toute équipe &
    Internationale & Cycle de vie complet \\
\norme{NIST SP 800-86} & IR intégrée & CISO, équipe SOC &
    USA + adopté monde & IR + forensique \\
\norme{RFC 3227} & Volatilité & Tout OPJ &
    Universelle & Ordre de collecte \\
ACPO Guide & Déontologie & Forces de l'ordre &
    Nationale + influence & Principes éthiques \\
SANS FOR508 & Technique avancée & Senior analyste &
    Formation privée & Forensique entreprise \\
\end{tableaucours}
\vspace{-0.8em}
\begin{center}{\small\itshape Tableau~7.6 --- Analyse DICES comparative des standards}
\end{center}

\subsection{Quel standard pour quelle situation~?}

\begin{tableaucours}{p{5.5cm}p{5.5cm}}{%
    \tableauHeader{Situation} &
    \tableauHeader{Standard(s) prioritaire(s)}
}
OPJ arrive sur une scène, matériel allumé &
    RFC~3227 (ordre de collecte) + ISO~27037 (procédure) \\
Analyse d'une image forensique en laboratoire &
    ISO~27041 (méthode) + ISO~27042 (analyse) \\
Investigation d'un incident ransomware en entreprise &
    NIST~SP~800-86 + SANS~FOR508 \\
Expert judiciaire devant un tribunal &
    ISO~27037 + ACPO Principe~3 + NIST~SP~800-86 (rapport) \\
Investigation impliquant un État partenaire Budapest &
    ISO~27043 + Convention Budapest art.~29 (conservation) \\
Formation d'un OPJ camerounais niveau débutant &
    ACPO 4~principes + RFC~3227 (ordre volatilité) \\
\end{tableaucours}
\vspace{-0.8em}
\begin{center}{\small\itshape Tableau~7.7 --- Sélection du standard selon la situation}
\end{center}

% ============================================================================
\section{Standards Émergents~: Cloud, IoT et Post-Quantique}
\label{sec:standards-emergents}
% ============================================================================

\subsection{Forensique cloud}

\begin{itemize}
    \item \norme{ISO/IEC~27050}~: \textit{Electronic Discovery}~;
          collecte de données électroniques pour procédures légales,
          incluant les environnements cloud.
    \item \norme{NIST SP~800-210}~: Cloud forensics challenges~;
          gestion des défis multitenancy, éphéméralité et multi-juridiction.
    \item \textbf{CSA Guidelines}~: Cloud Security Alliance,
          bonnes pratiques d'investigation dans les environnements cloud.
\end{itemize}

\subsection{Forensique IoT}

\begin{itemize}
    \item \norme{IEEE~P2933}~: Trusted IoT Data~;
          authenticité et intégrité des données produites par des objets connectés.
    \item \norme{ETSI TR~103~939}~: IoT testing methodology.
    \item \norme{ISO/IEC~30141}~: IoT reference architecture~;
          utile pour comprendre où trouver les données forensiques.
\end{itemize}

\subsection{Post-quantique et authenticité des preuves}

Le NIST finalise en 2024 les algorithmes post-quantiques de signature~:
\norme{CRYSTALS-Dilithium}, \norme{FALCON} et \norme{SPHINCS+}. Ces
algorithmes remplaceront à terme ECDSA et RSA pour la signature
cryptographique des preuves numériques. L'enjeu forensique~: les preuves
signées aujourd'hui avec ECDSA seront potentiellement cassables dans 10--15~ans.
Pour les affaires à longue durée (\textit{e.g.} terrorisme), la migration
vers les algorithmes post-quantiques est déjà une obligation éthique.

\begin{boxCRO}
\textbf{Analyse CRO --- Standards et choix méthodologique}

Le choix d'un standard n'est pas neutre vis-à-vis du Trilemme CRO~:

$\mathcal{R}$ (Fiabilité) est maximisée par les standards techniques~:
RFC~3227 (ordre de collecte), ISO~27037 (hash), ISO~27042 (méthodes validées).

$\mathcal{O}$ (Opposabilité) est maximisée par les standards à portée légale~:
ISO~27037 et ACPO sont reconnus par les tribunaux du monde entier~;
NIST~SP~800-86 est reconnu aux USA et dans de nombreux pays.

$\mathcal{C}$ (Confidentialité) est gérée par les standards procéduraux~:
ACPO Principe~1 (ne pas modifier) protège les données des parties~;
ISO~27041 (proportionnalité de la collecte) évite la sur-collecte.

\cro{$\uparrow$ ACPO Pr.1}{$\uparrow\uparrow$ ISO 27037/42}{$\uparrow\uparrow$ ISO 27037/ACPO}
\end{boxCRO}

\separateur

\begin{boxresume}
\textbf{Ce qu'il faut retenir de ce chapitre~:}
\begin{itemize}
    \item Les standards forensiques résolvent trois problèmes~:
          \textbf{interopérabilité} (preuves recevables partout),
          \textbf{opposabilité} (méthode défendable) et \textbf{formation}
          (base commune).

    \item \textbf{Pour le terrain~:} RFC~3227 (ordre de volatilité~:
          RAM $\to$ réseau $\to$ processus $\to$ disque) et ISO~27037
          (procédure de saisie, rôles DEFR/DES, hash SHA-256 minimum).

    \item \textbf{Pour l'analyse~:} ISO~27042 impose la séquence
          Préparation~$\to$ Extraction~$\to$ Analyse~$\to$ Interprétation~$\to$
          Rapport~; la distinction Examination/Analysis (NIST~SP~800-86)
          est la règle d'or du rapport forensique.

    \item \textbf{Pour la déontologie~:} ACPO 4~principes --- (1)~ne
          pas modifier~; (2)~compétence si modification~; (3)~audit trail~;
          (4)~responsabilité personnelle.

    \item \textbf{En droit camerounais~:} ISO~27037 est
          \textbf{substantiellement aligné} sur les articles~52--59 de la
          \loicam{2010/012}. La conformité à l'un implique la conformité
          à l'autre.

    \item \textbf{Choix situationnel~:} terrain~$\to$ RFC~3227 + ISO~27037~;
          laboratoire~$\to$ ISO~27041 + ISO~27042~; entreprise~$\to$ NIST + SANS~;
          tribunal~$\to$ ISO~27037 + ACPO~; transfrontalier~$\to$ Budapest art.~29.
\end{itemize}
\end{boxresume}

\bigskip

\begin{boxexo}[title={\ding{45}\quad Exercice~7.1 --- Sélection de standard \niveaubase}]
Pour chacune des situations suivantes, identifiez le(s) standard(s) à
appliquer en priorité et justifiez votre choix~:
\begin{enumerate}[(a)]
    \item Un OPJ camerounais arrive chez un suspect~; l'ordinateur est allumé,
          et un téléphone Android est posé sur le bureau en mode déverrouillé.
    \item Un analyste forensique doit déterminer si les 15~000~emails
          collectés dans une affaire de fraude ont été altérés depuis leur
          collecte.
    \item Une PME camerounaise victime d'un ransomware demande une investigation
          complète pour identifier l'attaquant et produire un rapport
          pour son assurance.
    \item Un expert judiciaire doit produire un rapport opposable dans
          une affaire impliquant des preuves stockées sur des serveurs
          Google (États-Unis).
\end{enumerate}
\end{boxexo}

\begin{boxexo}[title={\ding{45}\quad Exercice~7.2 --- Application ACPO + RFC~3227 \niveaumoyen}]
Vous intervenez dans l'affaire MVOM (Opération MVOM, Chapitre~\ref{chap:gr-affaires})~:
serveur Windows compromis par ransomware, allumé, connecté au réseau.
\begin{enumerate}[(a)]
    \item Appliquez l'ordre de volatilité RFC~3227 à cette scène spécifique~:
          listez les 6~premières actions dans l'ordre exact, en précisant
          l'outil et la commande correspondante.
    \item Pour chaque action, identifiez le principe ACPO auquel elle répond.
    \item L'analyste doit capturer la mémoire RAM (16~Go) sur un serveur
          en production. Le Principe~1 ACPO dit ``ne pas modifier les données''~;
          or la capture RAM modifie légèrement la RAM. Comment résoudre
          cette tension~? Quelle documentation apporter~?
    \item Rédigez l'entrée correspondante dans le formulaire de chaîne
          de custody (une ligne par action).
\end{enumerate}
\end{boxexo}

\begin{boxexo}[title={\ding{45}\quad Exercice~7.3 --- Analyse comparative DICES \niveauexpert}]
Un responsable forensique au sein de l'ANTIC doit rédiger une
\textit{note de politique forensique} recommandant les standards à adopter
officiellement au Cameroun.
\begin{enumerate}[(a)]
    \item Sur la base de l'analyse DICES, quel standard recommandez-vous
          comme fondement légal principal~? Justifiez en articulant portée
          légale, compatibilité avec la Loi~2010/012 et facilité de
          formation pour les OPJ.
    \item Identifiez deux tensions CRO que la mise en \oe uvre de ce standard
          créerait dans le contexte camerounais~; proposez une résolution
          pour chacune.
    \item La Convention de Budapest (art.~29) est complémentaire à ces
          standards~: expliquez en quoi, et dans quel cas un OPJ camerounais
          l'activerait en \emph{plus} d'ISO~27037.
\end{enumerate}
\end{boxexo}

\bigskip

\begin{boxinfo}
\textbf{Transition.} Le Chapitre~\ref{chap:custody} approfondit la
\textbf{chaîne de custody}~: le mécanisme procédural qui traduit concrètement
ISO~27037 et les principes ACPO en documents et formulaires opposables.
C'est le chapitre le plus important pour l'admissibilité des preuves
devant les tribunaux camerounais.
\end{boxinfo}